\documentclass[12pt, a4paper]{article}

% --- Packages ---
\usepackage[margin=1in]{geometry}
\usepackage{amsmath, amssymb, amsthm}
\usepackage{setspace}
\usepackage{hyperref}
\usepackage{natbib}
\usepackage{booktabs}
\usepackage{enumitem}
\usepackage{titlesec}

% --- Theorem environments ---
\newtheorem{proposition}{Proposition}
\newtheorem{lemma}{Lemma}
\newtheorem{remark}{Remark}
\newtheorem{definition}{Definition}

% --- Formatting ---
\onehalfspacing
\titleformat{\section}{\large\bfseries}{\thesection.}{0.5em}{}
\titleformat{\subsection}{\normalsize\bfseries}{\thesubsection.}{0.5em}{}

% --- Title ---
\title{\textbf{Referee Report: \\ ``Consumer Information and the Limits to Competition''} \\[6pt]
\large Mark Armstrong \& Jidong Zhou \\
\textit{American Economic Review}, 2022, 112(2), 534--577}
\author{Jordi Torres Vallverd\'{u} \\
\small Toulouse School of Economics --- MRes IO Theory, 2025--2026}
\date{\today}

\begin{document}

\maketitle
\thispagestyle{empty}

% =============================================
% PART 1: SUMMARY (2-4 pages)
% =============================================
\section{Summary}

\subsection{Motivation}

In many markets, firms supply differentiated products and compete in prices, while consumers possess incomplete information about how well each product matches their preferences. A central question for platforms (e.g., Amazon, Expedia, insurance comparison websites) and regulators is how to design the information environment that consumers face. More transparency improves product-consumer matching, but it also softens price competition by granting firms greater market power. Less transparency forces fierce price competition --- products are perceived as near-perfect substitutes --- but consumers may purchase a poorly matched product.

Armstrong and Zhou formalize this trade-off. They ask: within the class of all feasible signal structures, which one maximizes industry profit, and which one maximizes consumer surplus? The answer reveals that the interests of firms and consumers are sharply opposed in how information should be structured, yet both sides agree that information should be \emph{symmetric} across firms.

\subsection{Model Setup}

Consider an ex ante symmetric duopoly. Two firms ($i = 1,2$) each supply one variety of a differentiated product at zero marginal cost. A representative risk-neutral consumer values product $i$ at $v_i \geq 0$ and is initially uncertain about her valuations. She holds a symmetric prior over $(v_1, v_2)$, and only the \textbf{relative preference} $x \equiv v_1 - v_2$ matters for her choice.

\begin{itemize}[nosep]
    \item Let $F(x)$ denote the prior CDF of $x$, symmetric around $0$, with support $[-1,1]$.
    \item Let $\mu \equiv E[v_i]$ denote the expected valuation for a random product, and define the \textbf{differentiation parameter} $\delta \equiv E_F[\max\{x,0\}] = \int_{-1}^{0} F(x)\,dx$, which measures the expected match-quality gain from buying the preferred product over a random one.
    \item The full-information equilibrium price is $p_F = 1/(2f(0))$, where $f$ is the density of $F$.
\end{itemize}

Before purchasing, the consumer observes a private signal $s$ drawn from a signal structure $\{\sigma(s|x), S\}$, where $\sigma(s|x)$ specifies the distribution of signal $s$ when the true relative preference is $x$, and $S$ is the signal space. The signal structure is publicly known and chosen before firms set prices. After observing $s$, the consumer updates her belief about $x$ to a posterior distribution $F_s(x)$. Since she is risk-neutral, only the posterior mean $E_{F_s}[x]$ matters for her decision. Let $G(x)$ denote the CDF of this posterior mean across signals --- the \textbf{posterior distribution}.

Firms observe the posterior $G$ (since the signal structure is public) but not the consumer's realized signal $s$. They simultaneously set prices $p_1, p_2$ to maximize expected profit (Bertrand--Nash). The consumer buys from firm~1 if her expected relative preference satisfies $E[x \mid s] > p_1 - p_2$, and from firm~2 otherwise.

\textbf{Bayesian consistency} requires that $G$ be a \textbf{mean-preserving contraction} (MPC) of the prior $F$:
\begin{equation}
    \int_{-1}^{x} G(\tilde{x})\,d\tilde{x} \;\leq\; \int_{-1}^{x} F(\tilde{x})\,d\tilde{x}, \qquad \text{for all } x \in [-1,1], \text{ with equality at } x = 1.
\end{equation}
Any $G$ satisfying this constraint can be generated by some signal structure. This allows the authors to optimize directly over posteriors $G$ rather than over signal structures, which substantially simplifies the analysis.

\subsection{Bounds on Posteriors and the Role of Symmetry}

A key technical contribution is establishing bounds on the posterior distributions that can sustain a given pair of equilibrium prices $(p_1, p_2)$ as a pure-strategy Nash equilibrium of the pricing game.

\begin{lemma}[Lemma 1 in the paper]
A distribution $G$ implements prices $(p_1, p_2)$ as equilibrium prices if and only if
\begin{equation}
    L_{p_1,p_2}(x) \;\leq\; G(x) \;\leq\; U_{p_1,p_2}(x) \qquad \text{for all } x \in [-1,1],
\end{equation}
where the \textbf{lower bound} $L_{p_1,p_2}(x) = \max\{0,\; 1 - \pi_1/(p_2 + x)\}$ arises from firm~1's no-deviation condition, the \textbf{upper bound} $U_{p_1,p_2}(x) = \min\{\pi_2/(p_1 - x),\; 1\}$ arises from firm~2's no-deviation condition, and $\pi_i = p_i^2/(p_1 + p_2)$ is firm $i$'s equilibrium profit.
\end{lemma}

For symmetric prices $p_1 = p_2 = p$, these bounds simplify to:
\begin{equation}
    L_p(x) = \begin{cases} 0 & \text{if } x \leq -\frac{1}{2}p \\ 1 - \frac{p}{2(p+x)} & \text{otherwise} \end{cases}, \qquad
    U_p(x) = \begin{cases} \frac{p}{2(p-x)} & \text{if } x < \frac{1}{2}p \\ 1 & \text{otherwise}. \end{cases}
\end{equation}
The two bounds are mirror images: $L_p(x) = 1 - U_p(-x)$. A symmetric posterior $G$ can sustain price $p$ if and only if it lies between $L_p$ and $U_p$ and is an MPC of $F$.

A striking result is that symmetric signal structures dominate asymmetric ones for \emph{both} firms and consumers:

\begin{lemma}[Lemma 2 in the paper]
Relative to any asymmetric posterior that induces a pure-strategy equilibrium:
\begin{enumerate}[nosep]
    \item[(i)] there exists a symmetric posterior with \textbf{weakly higher profit for each firm};
    \item[(ii)] there exists a symmetric posterior with \textbf{weakly higher consumer surplus}.
\end{enumerate}
\end{lemma}

The intuition for part (i) is that when a signal favors one firm, the disfavored firm undercuts aggressively, triggering a price war that harms both firms --- even the favored one. For part (ii), asymmetric signals introduce both match inefficiency and price dispersion, and consumers are better off under some symmetric alternative.

\subsection{Main Results: Optimal Signal Structures}

\textbf{Firm-optimal signal.} The firm-optimal signal structure maximizes the sustainable symmetric price. The optimal posterior is $G = L_{p^*}$ for $x < 0$ (i.e., the posterior sits on the lower bound), yielding:
\begin{equation}
    p^* = \frac{2\delta}{1 - \log 2} \;\approx\; 6.5\,\delta.
\end{equation}
This posterior has a U-shaped density: mass is concentrated at the extremes ($x$ near $\pm 1$), meaning most consumers have strong preferences for one product. Few consumers are near-indifferent (near $x = 0$), so neither firm can profitably undercut. The consumer always buys her preferred product --- there is \textbf{no mismatch} --- and total welfare is maximized at $\mu + \delta$. Relative to full disclosure, this information structure raises industry profit by at least 63\% when the prior density is log-concave.

\textbf{Consumer-optimal signal.} The consumer-optimal signal structure pushes the posterior against the \emph{upper} bound $U_p$ for negative $x$. The optimal posterior is $G = \min\{F, U_{p^{**}}\}$ for $x < 0$, with:
\begin{equation}
    p^{**} = \frac{-\gamma}{1-\gamma}\, F^{-1}\!\left(\frac{\gamma}{2}\right), \qquad \text{where } \gamma \approx 0.05 \text{ solves } \gamma - \log \gamma = 3.
\end{equation}
This posterior amplifies the mass near $x = 0$: many consumers are nearly indifferent between products. Firms face intense price competition and set very low prices ($p^{**} \approx 0.05\, p_F$). However, near-indifferent consumers sometimes buy their less-preferred product, so there is \textbf{mismatch} and total welfare is not maximized. The price reduction more than compensates consumers for this match-quality loss.

\textbf{The limits to competition (Proposition 3).} Any industry profit $p \in [0, p^*]$ is feasible. For each profit level, the minimum consumer surplus is $\mu - \frac{1}{2}(1 + \log 2)\,p$ (attained by the lower-bound posterior $L_p$), and the maximum consumer surplus is $\mu + \delta - p$ when $p \geq p_F$ (efficient frontier), or a concave function of $p$ when $p < p_F$. A key implication of competition is that it is \textbf{impossible} to have both low profit and low consumer surplus simultaneously --- in contrast to the monopoly case studied by Roesler and Szentes (2017).


% =============================================
% PART 2: CRITIQUE (1-2 pages)
% =============================================
\section{Critique}

While the paper makes an important and technically elegant contribution, several aspects of the modeling framework and results deserve scrutiny.

\subsection{Scalar Heterogeneity and the Duopoly Restriction}

The paper reduces the consumer's two-dimensional preference $(v_1, v_2)$ to a scalar relative valuation $x = v_1 - v_2$. This is natural in the duopoly case, but it means that the \emph{level} of consumer willingness-to-pay (captured by $V = v_1 + v_2$) is effectively assumed away. In practice, both the level and the relative preference matter: a consumer who values both products highly behaves differently from one who values both poorly, even if their relative preference $x$ is identical. With more than two firms, preferences are inherently multidimensional, and the paper's approach --- reducing everything to a single scalar --- does not extend. The authors acknowledge this limitation and study three simple policies (``top product,'' ``top two,'' and full disclosure) for $n > 2$ firms, but the full characterization of optimal signal structures remains open.

\subsection{Pure-Strategy Restriction}

The focus on signal structures that induce pure-strategy pricing equilibria is motivated by tractability. The authors show (Proposition~4) that allowing mixed strategies improves consumer surplus by at most 1.6\%, which is reassuring. However, the firm-optimal signal under mixed strategies remains uncharacterized. More importantly, for some prior distributions and signal structures, pure-strategy equilibria may not exist at all, and the paper's framework provides no guidance for these cases. It would be valuable to understand whether the qualitative insights (e.g., symmetric signals dominate, the shape of the feasible frontier) extend when mixed strategies are the relevant equilibrium concept.

\subsection{Exogeneity of the Signal Structure}

The entire analysis assumes the signal structure is chosen exogenously by a designer and committed to before firms set prices. This ``centralized design'' assumption is strong. In many real markets, information is generated through a combination of firm advertising, consumer search, platform algorithms, and third-party reviews. Each of these has different strategic properties. In particular:
\begin{itemize}[nosep]
    \item If firms can influence the signal structure (e.g., through advertising), the optimal centralized policy may not be implementable.
    \item If consumers can choose how much information to acquire (endogenous search), the designer's ability to control the posterior is limited.
    \item The timing assumption --- signal structure chosen before prices --- is crucial. If firms could respond to the signal structure by adjusting product characteristics, the welfare conclusions could reverse.
\end{itemize}

\subsection{Symmetry of the Prior}

The paper assumes firms are ex ante symmetric: the prior over $(v_1, v_2)$ is symmetric between the two products. Many real-world markets feature vertical differentiation (one product is known to be better on average) or brand effects. The authors suggest (Section~IV) that extending to asymmetric priors is an interesting direction, but they do not explore whether their key results --- particularly the optimality of symmetric signals (Lemma~2) --- survive. With an asymmetric prior, biased signals that further favor the better firm might no longer be dominated, and the structure of the feasible set could change qualitatively.

\subsection{Welfare Metric and the Role of Mismatch}

The consumer-optimal signal structure deliberately introduces mismatch: some consumers buy their less-preferred product. The paper treats this as a pure efficiency loss, measured by the reduction in $\int_{-1}^{0} G(x)\,dx$ relative to full information. However, in practice, product mismatch may have dynamic consequences --- consumer regret, returns, negative reviews --- that go beyond the static welfare calculation. If mismatch generates costs not captured in the model, the consumer-optimal policy may be less attractive than the static analysis suggests.


% =============================================
% PART 3: EXTENSION (2-4 pages)
% =============================================
\section{Extension: Information Designer as a Strategic Agent}

The paper explicitly calls for future work embedding the information designer as a strategic economic agent (Section~IV, p.~562). We take up this suggestion by introducing a platform that designs the information environment and extracts rents through fees, and we investigate how its welfare implications differ from the benchmark cases.

\subsection{Setup}

We augment the Armstrong--Zhou framework with a profit-maximizing platform that controls the signal structure and charges a per-transaction fee $\tau \geq 0$ to firms. The platform commits to a signal structure $G$ (equivalently, a posterior distribution that is an MPC of $F$) and a fee $\tau$ before firms simultaneously set prices. The timing is:

\begin{enumerate}[nosep]
    \item The platform publicly commits to a symmetric signal structure (inducing posterior $G$) and a per-unit fee $\tau$ charged to each firm on every sale.
    \item Firms observe $(G, \tau)$ and simultaneously choose prices $p_1, p_2$ to maximize their own profit net of the fee.
    \item The consumer observes a private signal, updates beliefs, observes prices, and purchases.
\end{enumerate}

We maintain the duopoly setup: two symmetric firms, zero marginal production cost, prior $F$ with support $[-1,1]$, and the consumer always buys from one of the two firms (outside option irrelevant). Each firm's effective marginal cost is now $\tau$ rather than zero.

\subsection{Pricing Equilibrium with Fees}

In the pricing subgame, each firm maximizes $(p_i - \tau) \cdot q_i(p_1, p_2; G)$, where $q_i$ is firm $i$'s demand. By a standard change of variables, define $\tilde{p}_i \equiv p_i - \tau$ as the firm's markup over cost. The pricing game in terms of markups $(\tilde{p}_1, \tilde{p}_2)$ is identical to the original Armstrong--Zhou game with posterior $G$ and zero cost. Therefore, the equilibrium markup is determined entirely by $G$, exactly as before.

\begin{proposition}[Pricing with fees]
If the posterior $G$ induces a symmetric pure-strategy equilibrium with price $p$ in the original (zero-cost) game, then the same $G$ induces a symmetric equilibrium with price $p + \tau$ in the game with per-unit fee $\tau$. Each firm's equilibrium markup is $\tilde{p} = p$, and each firm's profit is $\pi^{\text{firm}} = p/2$, independent of $\tau$. The consumer faces price $p + \tau$.
\end{proposition}

This cost pass-through result follows from the structure of the Bertrand game with differentiated products: the no-deviation conditions depend on the markup, not the price level. Importantly, the fee $\tau$ is fully passed through to the consumer.

\subsection{Platform's Problem}

The platform's revenue is $\tau$ per transaction (total revenue $= \tau$ since one unit is always sold). The platform chooses $(G, \tau)$ to maximize:
\begin{equation}
    \Pi^{\text{plat}}(\tau) = \tau,
\end{equation}
subject to:
\begin{enumerate}[nosep]
    \item $G$ is an MPC of $F$ (Bayesian consistency);
    \item $G$ supports a symmetric pure-strategy equilibrium with some markup $\tilde{p}(G)$;
    \item Participation: each firm earns non-negative profit, i.e., $\tilde{p}(G) \geq 0$ (always satisfied);
    \item (Optional) A constraint on consumer participation or welfare.
\end{enumerate}

If the platform faces no constraint beyond equilibrium existence, it wants $\tau$ as large as possible. Since $\tau$ is simply added to the consumer's price and does not affect firm behavior, the platform is unconstrained in $\tau$ as long as the consumer still prefers buying from a firm over the outside option. When the outside option is irrelevant ($V > 2$), any $\tau$ is feasible, and the platform's problem is trivially unbounded.

This motivates introducing a meaningful constraint. We consider two cases.

\subsubsection{Case 1: Consumer Participation Constraint}

Suppose the consumer has an outside option yielding utility zero. The consumer buys product $i$ only if $v_i - (p_i) \geq 0$, i.e., her expected utility from buying must be non-negative. In a symmetric equilibrium with price $p + \tau$ and posterior $G$, the consumer's expected surplus is:
\begin{equation}
    CS(G, \tau) = \mu + \int_{-1}^{0} G(x)\,dx - (p(G) + \tau),
\end{equation}
where $p(G)$ denotes the equilibrium markup induced by $G$. The participation constraint is $CS \geq 0$, giving:
\begin{equation}
    \tau \;\leq\; \mu + \int_{-1}^{0} G(x)\,dx - p(G) \;\equiv\; \bar{\tau}(G).
\end{equation}

The platform maximizes $\tau = \bar{\tau}(G)$ over all feasible $G$. This is equivalent to maximizing the consumer's \emph{gross} surplus $\mu + \int_{-1}^0 G(x)\,dx - p(G)$ and then extracting it entirely via $\tau$.

\begin{proposition}[Platform-optimal signal with participation constraint]
The platform-optimal signal structure maximizes $\int_{-1}^{0} G(x)\,dx - p(G)$, which is the consumer's match-quality gain net of the markup. The platform faces the same trade-off as in the consumer-optimal problem: lowering $p(G)$ requires concentrating $G$ near $x = 0$ (more indifferent consumers), which reduces $\int_{-1}^0 G(x)\,dx$ (more mismatch). The optimal $G$ solves:
\begin{equation}
    \max_{G, p} \quad \int_{-1}^{0} G(x)\,dx \;-\; p \qquad \text{s.t.} \quad L_p(x) \leq G(x) \leq U_p(x), \quad G \text{ is MPC of } F.
\end{equation}
This is precisely the consumer-optimal problem from Armstrong and Zhou. Therefore, the platform-optimal signal \textbf{coincides} with the consumer-optimal signal: $G^{\text{plat}} = G^{**}$, $p^{\text{plat}} = p^{**}$, and the platform extracts all consumer surplus via $\tau^* = CS(G^{**}, 0)$.
\end{proposition}

\begin{remark}
This result echoes a standard insight: a monopoly intermediary that can set access fees will choose the downstream market structure that maximizes total downstream surplus and then extract it. Here, the ``downstream surplus'' available to the platform is consumer surplus (since firm profits cannot be extracted via $\tau$ --- markups are invariant to $\tau$). Thus the platform designs information to maximize consumer surplus and appropriates it.
\end{remark}

\textbf{Welfare implications.} Under the platform-optimal policy:
\begin{itemize}[nosep]
    \item Consumers get exactly zero surplus (participation binds).
    \item Firms earn $p^{**}/2 \approx 0.025\, p_F$ each --- low but positive.
    \item The platform captures $\tau^* = \mu + \int_{-1}^0 G^{**}(x)\,dx - p^{**}$, which is the bulk of welfare.
    \item Total welfare is $\mu + \int_{-1}^0 G^{**}(x)\,dx < \mu + \delta$ (mismatch occurs).
\end{itemize}
Relative to the Armstrong--Zhou benchmark without a platform, consumers are strictly worse off (zero surplus vs.\ positive surplus under the consumer-optimal signal), and total welfare is lower than under the firm-optimal signal.

\subsubsection{Case 2: Platform Charges Firms a Lump-Sum Fee}

Now suppose the platform charges each firm a symmetric lump-sum participation fee $\phi$ to be listed, rather than a per-unit fee. Firms decide whether to participate before prices are set. The timing becomes:

\begin{enumerate}[nosep]
    \item Platform commits to $(G, \phi)$.
    \item Each firm decides whether to pay $\phi$ and participate.
    \item Participating firms set prices; consumer buys.
\end{enumerate}

In a symmetric equilibrium, both firms participate if and only if their equilibrium profit covers the fee: $\pi^{\text{firm}}(G) = p(G)/2 \geq \phi$. The platform's revenue is $2\phi$, so it maximizes:
\begin{equation}
    \Pi^{\text{plat}} = 2\phi = p(G),
\end{equation}
subject to $G$ being an MPC of $F$ and inducing a pure-strategy equilibrium.

\begin{proposition}[Platform-optimal signal with lump-sum fees]
When the platform extracts rents via lump-sum fees, the platform-optimal signal structure \textbf{coincides with the firm-optimal signal}: $G^{\text{plat}} = L_{p^*}$ with $p^* = 2\delta/(1 - \log 2)$. The platform extracts $\Pi^{\text{plat}} = p^* \approx 6.5\,\delta$, consumers buy their preferred product (no mismatch), and total welfare is maximized at $\mu + \delta$. Consumer surplus is $\mu + \delta - p^* - 0 = \mu + \delta - p^*$ (since the lump-sum fee does not affect prices).
\end{proposition}

\textbf{Welfare implications.} Under lump-sum fees, the platform's incentives are fully aligned with firms: it wants to maximize industry profit so it can extract it. This leads to the firm-optimal signal, which also maximizes total welfare. Consumers face high prices ($p^*$) but buy their preferred product. The comparison between the two fee structures is stark:

\begin{center}
\renewcommand{\arraystretch}{1.3}
\begin{tabular}{l c c}
\toprule
 & \textbf{Per-unit fee $\tau$} & \textbf{Lump-sum fee $\phi$} \\
\midrule
Signal structure & Consumer-optimal ($G^{**}$) & Firm-optimal ($L_{p^*}$) \\
Equilibrium price & $p^{**} + \tau^*$ (high) & $p^*$ (high) \\
Mismatch & Yes & No \\
Total welfare & $< \mu + \delta$ & $= \mu + \delta$ (max) \\
Consumer surplus & $0$ & $\mu + \delta - p^* > 0$ \\
Platform revenue & $\tau^* = CS(G^{**},0)$ & $p^*$ \\
\bottomrule
\end{tabular}
\end{center}

\subsection{Discussion: Does the Platform Prefer Higher or Lower Prices?}

This extension reveals that the answer to the question ``does the information designer prefer higher or lower downstream prices?'' depends on how the designer monetizes.

\begin{itemize}
    \item \textbf{Per-unit fees} (commission model, e.g., Expedia, Uber): the platform wants to maximize the \emph{consumer's willingness to pay} net of the firm's markup, so it designs information to minimize the markup. This aligns the platform with consumers in terms of signal design (low prices, high competition), but the platform then extracts the consumer surplus itself. Consumers end up with nothing, and mismatch occurs.

    \item \textbf{Lump-sum fees} (listing/subscription model, e.g., premium marketplace listings): the platform wants to maximize \emph{firm profits}, so it designs information to soften competition. This produces the firm-optimal signal with no mismatch, maximizing total welfare. Paradoxically, consumers are better off under this extractive platform than under the per-unit-fee platform, because they retain $\mu + \delta - p^*$ in surplus (which is positive as long as $\mu > p^* - \delta$, i.e., the basic product utility is large enough).
\end{itemize}

This suggests that from a \textbf{regulatory perspective}, the fee structure of platforms is not merely a distributional issue --- it affects the platform's incentives to design information, which in turn affects allocative efficiency. A regulator concerned with total welfare might prefer platforms that use lump-sum fees, even though such platforms appear more ``extractive'' toward firms, because they align information design with efficiency.

\subsection{Sketch: Competing Platforms}

A natural next step is to consider two competing platforms, each designing its own information environment and charging fees. If consumers can choose which platform to use, platforms compete to attract consumers, which constrains their ability to extract surplus. In the limit of intense platform competition, the equilibrium information design should approach the consumer-optimal signal, as platforms use favorable information design as a competitive tool. This connects to the broader literature on platform competition (e.g., Rochet and Tirole, 2003) and suggests that market structure at the platform level --- not just at the product level --- is a key determinant of information provision.

Formalizing this requires specifying how consumers choose between platforms and whether firms can multi-home, which we leave for future work.

% =============================================
% REFERENCES
% =============================================
\section*{References}

\begin{itemize}[nosep, leftmargin=*]
    \item Armstrong, M. and J. Zhou (2022). ``Consumer Information and the Limits to Competition.'' \textit{American Economic Review}, 112(2), 534--577.
    \item Kamenica, E. and M. Gentzkow (2011). ``Bayesian Persuasion.'' \textit{American Economic Review}, 101(6), 2590--2615.
    \item Roesler, A.-K. and B. Szentes (2017). ``Buyer-Optimal Learning and Monopoly Pricing.'' \textit{American Economic Review}, 107(7), 2072--2080.
    \item Anderson, S. and R. Renault (2006). ``Advertising Content.'' \textit{American Economic Review}, 96(1), 93--113.
    \item Johnson, J. and D. Myatt (2006). ``On the Simple Economics of Advertising, Marketing, and Product Design.'' \textit{American Economic Review}, 96(3), 756--784.
    \item Bergemann, D., B. Brooks, and S. Morris (2015). ``The Limits of Price Discrimination.'' \textit{American Economic Review}, 105(3), 921--957.
    \item Rochet, J.-C. and J. Tirole (2003). ``Platform Competition in Two-Sided Markets.'' \textit{Journal of the European Economic Association}, 1(4), 990--1029.
\end{itemize}

\end{document}
